\documentclass{article}
\usepackage[UTF8]{ctex}
\usepackage{amsmath,amssymb,amsfonts}
\usepackage{enumitem}
\usepackage{geometry}
\geometry{a4paper, margin=1in}
% === 错题本排版优化（不参与题目标识） ===
\setlength{\parindent}{0pt}
\setlength{\parskip}{0.35em plus 0.10em minus 0.08em}
\setlength{\emergencystretch}{3em}
\setlist{itemsep=0.25em, topsep=0.40em, parsep=0pt}
\allowdisplaybreaks[2]
% === 错题本排版优化结束 ===
\begin{document}
\section*{中国矿业大学(北京)2026年数学分析真题}
\section*{27．中国矿业大学（北京）2026年数学分析试卷（回忆版）}

\subsection{题目1}
1． \(\lim _{\alpha \rightarrow 0} \int_{\alpha}^{1+\alpha} \frac{\mathrm{d} x}{1+x^{2}+\alpha^{2}}=\) \(\_\_\_\_\)．

\subsection{题目2}
2．\(\sum_{n=1}^{\infty}(-1)^{n} \frac{2 n-1}{(2 n+3)!}=\) \(\_\_\_\_\)．

\subsection{题目3}
3． \(\int \frac{\ln x}{(1-x)^{2}} \mathrm{~d} x=\) \(\_\_\_\_\)．

\subsection{题目4}
4．高阶无穷小问题

\subsection{题目5}
5．\(f(x)=\frac{1}{x^{2}} \arctan \frac{1}{x-\frac{1}{x}}\) 的第一类间断点

\subsection{题目6}
6．\(x=a \cos t, y=a \sin t, z=a t\) ，求曲线积分 \(\int_{L} \frac{z^{2}}{x^{2}+y^{2}} \mathrm{~d} s=\) \(\_\_\_\_\)．

\subsection{题目7}
7． \(\int_{0}^{2 \pi} \mathrm{~d} x \int_{x}^{2 \pi} \frac{\sin ^{2} y}{y^{2}} \mathrm{~d} x=\) \(\_\_\_\_\)．

\subsection{题目8}
8．反常积分 \(\int_{0}^{+\infty} \frac{\sin x}{x} \mathrm{~d} x=\) \(\_\_\_\_\)（填＂绝对收敛＂或＂条件收敛＂或＂发散＂）。

\subsection{题目9}
9．\(y=f(x)\) 与 \(y=x^{2}-x\) 在 \((1,0)\) 点有公共的切线，求 \(\lim _{n \rightarrow \infty} f\left(\frac{n}{n+2}\right)=\) \(\_\_\_\_\)．

\subsection{题目10}
10．求 \(y=a x, y=b x, x+y=c, x+y=d\) 所围曲面的面积 \(\_\_\_\_\)

二．证明 \(f(x)=e^{x} \cos \frac{1}{x}\) 在 \((0,1)\) 上的一致连续性．

三．求 \(f(x)=\frac{2 x^{2}}{(1-x)^{2}}\) 的极值点与极值，凹凸区间，拐点与渐近线．

四．求曲面积分

\[
\iint_{\Sigma} \frac{x^{2} y \mathrm{~d} y \mathrm{~d} z+y^{2} z \mathrm{~d} z \mathrm{~d} x+\left(z^{2}+1\right) \mathrm{d} x \mathrm{~d} y}{4 x^{2}+4 y^{2}+z^{2}} .
\]

其中 \(\sum: 4 x^{2}+4 y^{2}+z^{2}=4 ; z \geq 0\) ．

五．叙述有限覆盖定理与连续函数的根的存在性定理，并用有限覆盖定理证明根的存在性定理．

六．证明 \(f_{n}(x)=n x(1-x)^{n}\) 在 \([0,1]\) 上的一致收敛性．

\begin{itemize}
  \item[七．] 设 \(a_{1}=3, a_{n+1}=\frac{1}{2} a_{n}+\frac{1}{a_{n}}\) ．
  \item[（1）] 证明：\(\left\{a_{n}\right\}\) 收敛，并求其极限．
  \item[（2）] 证明：对任意正数 \(p\) ，数项级数 \(\sum_{n=1}^{\infty} n^{p}\left(\frac{a_{n}}{a_{n+1}}-1\right)\) 收敛。
\end{itemize}

八．\(f(x)\) 在 \([0,1]\) 上连续，\(I=\int_{0}^{1} f(x) \mathrm{d} x \neq 0\) ，证明：在 \([0,1]\) 内存在两个不同的点 \(x_{1}, x_{2}\) ，使得

\[
\frac{1}{f\left(x_{1}\right)}+\frac{1}{f\left(x_{2}\right)}=\frac{2}{I} .
\]

九．求 \(f(x, y)=3(x-2 y)^{2}+x^{3}-8 y^{3}\) 的极值．
\end{document}
